\documentclass[11pt,a4paper]{article}

\usepackage[utf8]{inputenc}
\usepackage[T1]{fontenc}
\usepackage{lmodern}
\usepackage{geometry}
\geometry{a4paper, margin=1in}

\usepackage{mathpazo} % Palatino font
\usepackage{titlesec}
\usepackage{hyperref}
\usepackage{graphicx}
\usepackage{xcolor}
\usepackage{listings}
\usepackage{enumitem}

% Setup links
\hypersetup{
    colorlinks=true,
    linkcolor=blue,
    filecolor=magenta,      
    urlcolor=cyan,
    pdftitle={The Complexity of Credit Risk Data Engineering},
}

% Setup code blocks
\definecolor{codegray}{rgb}{0.5,0.5,0.5}
\definecolor{codepurple}{rgb}{0.58,0,0.82}
\definecolor{backcolour}{rgb}{0.95,0.95,0.92}

\lstdefinestyle{mystyle}{
    backgroundcolor=\color{backcolour},   
    commentstyle=\color{codegray},
    keywordstyle=\color{blue}\bfseries,
    numberstyle=\tiny\color{codegray},
    stringstyle=\color{codepurple},
    basicstyle=\ttfamily\footnotesize,
    breakatwhitespace=false,         
    breaklines=true,                 
    captionpos=b,                    
    keepspaces=true,                 
    numbers=none,                    
    showspaces=false,                
    showstringspaces=false,
    showtabs=false,                  
    tabsize=2
}

\lstset{style=mystyle}

\title{\Huge \textbf{The Engineering Reality of Credit Risk Data}}
\author{Architecture Review: Goto Financial / Home Credit Pipeline}
\date{\today}

\begin{document}

\maketitle

\section*{Executive Summary}
As highlighted by the \textbf{1st Place Kaggle Winner} of the Home Credit Default Risk competition, credit underwriting is one of the most complex domains for applying machine learning. The data is fundamentally \textit{heterogeneous}, collected over different timelines, and originates from diverse downstream systems. 

While the theoretical concepts of ``smart features'' are well understood, this paper explores exactly how our proprietary codebase (\texttt{src/hc\_risk/features.py}) implements these principles to mathematically sanitize and compress volatile banking data.

\section{Data Cleansing \& Safety Architecture}
Before any predictive modeling can occur, the raw relational data must be made safe for consumption by gradient boosting algorithms. Left unchecked, silent data anomalies in finance completely invert risk profiles.

\subsection{1. Sentinel Value Decoupling}
In legacy systems, missing data is rarely recorded as \texttt{NULL}. Frequently, programmers inject ``sentinel values'' (impossible extreme numbers) as placeholders. Our codebase handles the most notorious of these—the employment tenure sentinel:

\begin{lstlisting}[language=Python]
# Handling the 'Days Employed' sentinel
frame["DAYS_EMPLOYED_ANOM"] = (frame["DAYS_EMPLOYED"] == 365243).astype(int)
frame.loc[frame["DAYS_EMPLOYED"] == 365243, "DAYS_EMPLOYED"] = np.nan
\end{lstlisting}

A value of 365,243 equates to 1,000 years of employment. If passed blindly into a machine learning model, the algorithm concludes that these applicants possess infinite job stability, drastically mispricing their risk. By introducing an explicit \texttt{DAYS\_EMPLOYED\_ANOM} binary flag and nullifying the raw value, the model correctly isolates the ``unemployed/retired'' segment as a distinct risk cohort.

\subsection{2. Ratio Engineering Constraints}
The backbone of the \textit{Capacity Pillar} involves dividing debt by income. In a large-scale database, dividing by zero occurs almost constantly (e.g., applicants with \$0 reported debt). A single \texttt{ZeroDivisionError} will cascade and fail an automated pipeline. 

Our explicit constraint wrapper resolves this:
\begin{lstlisting}[language=Python]
def safe_divide(numerator: pd.Series, denominator: pd.Series):
    if isinstance(denominator, pd.Series):
        denominator = denominator.replace({0: np.nan})
    return numerator / denominator
\end{lstlisting}

This transforms a fatal programming error into a \texttt{NaN} (Not a Number) value. Tree-based learners natively route \texttt{NaN} values down designated split paths, converting a mathematical exception into a valid predictive signal (i.e., ``the borrower has no history of this denominator'').

\subsection{3. Memory Optimization for Flattened Matrices}
The process of flattening 1-to-N transaction logs into single applicant profiles creates massively wide data matrices. Our feature generation pipeline outputs upwards of 700 to 1,000 distinct columns. To prevent out-of-memory (OOM) crashes in production, the codebase aggressively downcasts 64-bit precision formats immediately upon ingestion:

\begin{lstlisting}[language=Python]
def reduce_memory(df: pd.DataFrame) -> pd.DataFrame:
    for column in df.columns:
        if pd.api.types.is_float_dtype(df[column]):
            df[column] = pd.to_numeric(df[column], downcast="float")
        elif pd.api.types.is_integer_dtype(df[column]):
            df[column] = pd.to_numeric(df[column], downcast="integer")
    return df
\end{lstlisting}
This deterministic downcasting shrinks the dataframe's RAM footprint by 40-50\%, enabling in-memory joins during distributed model training.

\section{The ``Many-to-One'' Aggregation Pipeline}
The winner explicitly notes: \textit{``These features employed various forms of aggregation on many-to-one tables, and joined them with application features.''} 

Our codebase implements this via distinct temporal aggregation functions designed to compress transaction histories into behavioral summaries.

\subsection{Recency and Time-Decay Targeting}
Historical data loses predictive power as it ages. The codebase prevents diluted aggregates via intentional subsetting:
\begin{itemize}
    \item \textbf{\texttt{\_top\_k\_aggregate}}: Filters historical transaction tables (such as \texttt{installments\_payments.csv}) down to a sliding window of the most recent $K$ interactions before averaging default rates.
    \item \textbf{\texttt{\_lag\_numeric\_features}}: Constructs explicit time-series deltas, measuring velocity (e.g., has the borrower's payment behavior improved or degraded between Month -6 and Month -1?).
    \item \textbf{\texttt{\_latest\_record\_features}}: Discards the entire history and locks onto the absolute final interaction state (e.g., is the credit card currently active or closed?).
\end{itemize}

\section{Conclusion}
While modern algorithms (\texttt{LightGBM}, \texttt{XGBoost}) provide the statistical heavy lifting, they are entirely dependent on the structural integrity of the input data. Our codebase demonstrates that applying machine learning to credit underwriting is less about algorithmic tuning and vastly more dependent on robust, defensively engineered data pipelines that mathematically capture human financial behavior.

\end{document}
